\subsection{Sets, Maps, Relations, and Universal Properties}

\begin{theorem} \label{thm:quotient-universal-property} \hlt{Universal Property of Quotients}\\
Let $\pi:\mathcal{V}\to\mathcal{V}/\mathcal{U}$ be the quotient map.\\
If $L\in\LL(\mathcal{V},\mathcal{W})$ and $\mathcal{U}\subseteq\ker(L)$, then $\exists!\,\overline{L}\in\LL(\mathcal{V}/\mathcal{U},\mathcal{W})$ such that $\overline{L}\pi=L$, $\overline{L}(\mathbf{v}+\mathcal{U})=L\mathbf{v}$.\\
Moreover, $\ker(\overline{L})=\ker(L)/\mathcal{U}$ and $\operatorname{range}(\overline{L})=\operatorname{range}(L)$.
\begin{proof}
If $\mathbf{v}+\mathcal{U}=\mathbf{w}+\mathcal{U}$, then $\mathbf{v}-\mathbf{w}\in\mathcal{U}\subseteq\ker(L)$, so $L\mathbf{v}=L\mathbf{w}$. Hence the formula defines a linear map $\overline{L}$ and gives $\overline{L}\pi=L$. If $M\pi=L$, then $M(\mathbf{v}+\mathcal{U})=M\pi(\mathbf{v})=L\mathbf{v}=\overline{L}(\mathbf{v}+\mathcal{U})$, so $M=\overline{L}$.\\
Finally, $\mathbf{v}+\mathcal{U}\in\ker(\overline{L}) \iff L\mathbf{v}=\mathbf{0} \iff \mathbf{v}\in\ker(L)$, giving $\ker(\overline{L})=\ker(L)/\mathcal{U}$.\\
The range equality follows from $\overline{L}(\mathbf{v}+\mathcal{U})=L\mathbf{v}$.\\
\hlt{Structural Commentary}: Maps killing $\mathcal{U}$ descend uniquely to $\mathcal{V}/\mathcal{U}$.\\
\hlt{Key Insight}: Quotients are the natural domain for maps whose ambiguity contains $\mathcal{U}$.
\end{proof}
\end{theorem}

\begin{figure}[H]
  \centering
  \begin{tikzpicture}[
  x=0.88cm,
  y=0.88cm,
  >=Stealth,
  every node/.style={font=\small},
  point/.style={circle, fill=gray!75, inner sep=1.55pt},
  arrow/.style={->, gray!65!black, line width=0.7pt},
  factor/.style={->, dashed, gray!65!black, line width=0.7pt},
  hint/.style={font=\scriptsize, text=gray!70!black, inner sep=1pt}
]
  \node[point] (V) at (0,0.75) {};
  \node[point] (Q) at (0,-0.75) {};
  \node[point] (W) at (3.1,0.75) {};

  \draw[arrow] (V) -- node[above, hint] {$L$} (W);
  \draw[arrow] (V) -- node[left, hint] {$\pi$} (Q);
  \draw[factor] (Q) -- node[above, sloped, hint] {$\overline{L}$} (W);

  \node at (0,1.08) {$\mathcal{V}$};
  \node at (3.1,1.08) {$\mathcal{W}$};
  \node at (0,-1.1) {$\mathcal{V}/\mathcal{U}$};
\end{tikzpicture}

  \caption{If $\mathcal{U}\subseteq\ker(L)$, then $L=\overline{L}\pi$ uniquely.}
  \label{fig:quotient-universal-property}
\end{figure}
